\section*{Capítulo \ref{ch:g11:becoming-male-female} --- Llegar a ser varón o mujer}

\begin{solution}{exo:g11:becoming-male-female:1}
El del padre: todo óvulo lleva un X, mientras que la mitad de los
espermatozoides llevan un X y la mitad un Y, así que XX y XY son
igualmente probables.
\end{solution}

\begin{solution}{exo:g11:becoming-male-female:2}
Dos gónadas indiferenciadas, cada una con dos conductos al lado: un
conducto de Wolff (futura vía masculina) y un conducto de Müller
(futura vía femenina). Idénticos en los embriones XX y XY.
\end{solution}

\begin{solution}{exo:g11:becoming-male-female:3}
En el \omterm{def:g11:cell-cycle-mitosis:chromosome}{cromosoma} Y; se expresa en la gónada indiferenciada hacia la
semana 7, y su \omterm{def:g11:gene-expression:protein}{proteína} enciende los \omterm{def:g10:universal-dna:gene}{genes} que convierten la gónada en
un testículo.
\end{solution}

\begin{solution}{exo:g11:becoming-male-female:4}
La \omterm{prop:g11:becoming-male-female:hormones}{testosterona} mantiene los conductos de Wolff y los desarrolla hasta
formar la vía masculina; la hormona antimülleriana hace regresar los
conductos de Müller.
\end{solution}

\begin{solution}{exo:g11:becoming-male-female:5}
Un centro del encéfalo libera pulsos de una hormona que hace que la
hipófisis segregue LH y FSH, que activan las gónadas. Caracteres
sexuales secundarios: los cambios que producen las hormonas sexuales
--- vello, voz, mamas, reparto del músculo y de la grasa, el estirón
del crecimiento y el crecimiento de los órganos reproductores.
\end{solution}

\begin{solution}{exo:g11:becoming-male-female:6}
Desarrolló una vía femenina a pesar de su \omterm{def:g11:cell-cycle-mitosis:chromosome}{cromosoma} Y: la vía masculina
exige las hormonas del testículo, y la femenina se forma sin ninguna
gónada.
\end{solution}

\begin{solution}{exo:g11:becoming-male-female:7}
La \omterm{prop:g11:becoming-male-female:hormones}{testosterona} construye los conductos masculinos pero no suprime los
femeninos; para eso el testículo debe segregar una segunda sustancia
--- la AMH.
\end{solution}

\begin{solution}{exo:g11:becoming-male-female:8}
Testículos (el Sry actúa); conductos de Wolff conservados y los de
Müller regresados; anatomía externa masculina. Estéril: producir
espermatozoides exige \omterm{def:g10:universal-dna:gene}{genes} del Y que el transgén no aporta.
\end{solution}

\begin{solution}{exo:g11:becoming-male-female:9}
Estrógenos: el 50\% hacia los 11 y la meseta hacia los 14 o 15.
\omterm{prop:g11:becoming-male-female:hormones}{Testosterona}: el 50\% hacia los 13 y la meseta hacia los 16 o 17.
\end{solution}

\begin{solution}{exo:g11:becoming-male-female:10}
Los huesos dejan de crecer cuando la hormona sexual llega a su nivel
adulto; los estrógenos lo alcanzan unos dos años antes que la
\omterm{prop:g11:becoming-male-female:hormones}{testosterona}, así que los \omterm{def:g10:muscles-and-joints:joint}{cartílagos} de crecimiento de las chicas se
cierran antes.
\end{solution}

\begin{solution}{exo:g11:becoming-male-female:11}
Sin testículo: la gónada se convierte en una gónada de tipo ovárico;
sin \omterm{prop:g11:becoming-male-female:hormones}{testosterona} ni AMH: los conductos de Müller se desarrollan, los de
Wolff regresan y la anatomía externa es femenina al nacer. En la
\omterm{prop:g11:becoming-male-female:puberty}{pubertad}, las gónadas, sin óvulos normales, segregan poco; sin
tratamiento hormonal los caracteres secundarios se desarrollan mal.
\end{solution}

\begin{solution}{exo:g11:becoming-male-female:12}
Testículos (el \omterm{prop:g11:becoming-male-female:sry}{SRY} funciona). Los conductos de Wolff regresan (no hay
respuesta a la \omterm{prop:g11:becoming-male-female:hormones}{testosterona}); los de Müller regresan también (la AMH
actúa): no hay vía interna de ninguno de los dos tipos. Anatomía
externa femenina (su forma masculina necesita la acción de la
\omterm{prop:g11:becoming-male-female:hormones}{testosterona}). En la \omterm{prop:g11:becoming-male-female:puberty}{pubertad}, la \omterm{prop:g11:becoming-male-female:hormones}{testosterona} convertida en estrógenos
produce mamas y una constitución femenina, sin vello corporal.
\end{solution}

\begin{solution}{exo:g11:becoming-male-female:13}
Una vez que una gónada se ha convertido en testículo o en ovario, sus
hormonas podrían modelar los conductos y el cuerpo exactamente igual
que en los mamíferos; lo que difiere es el primer paso --- un proceso
sensible a la temperatura sustituye al \omterm{prop:g11:becoming-male-female:sry}{SRY} como interruptor que decide
la gónada.
\end{solution}

\begin{solution}{exo:g11:becoming-male-female:14}
El X lleva cientos de \omterm{def:g10:universal-dna:gene}{genes} que necesitan todas las \omterm{def:g10:cells-common-unit:cell}{células}, así que
ningún embrión sobrevive sin uno; el Y no lleva casi nada aparte de los
\omterm{def:g10:universal-dna:gene}{genes} que determinan el testículo y los relacionados con los
espermatozoides, así que a las personas XX no les falta nada que
necesiten. Un Y solo pasa a los hijos varones porque las hijas reciben
el X del padre.
\end{solution}

\begin{solution}{exo:g11:becoming-male-female:15}
Significa que la vía femenina se desarrolla cuando no interviene
ninguna hormona testicular --- un hecho sobre qué señales hacen falta,
no una clasificación de importancia. Refleja un plan embrionario común
sobre el que un linaje de \omterm{def:g10:cells-common-unit:cell}{células}, el testículo, impone una alternativa
mediante dos hormonas: los dos sexos son dos resultados de un mismo
programa, y se diferencian por un interruptor.
\end{solution}

\begin{solution}{pb:g11:becoming-male-female:1}
\textbf{1.} El ovario no hace falta para la vía femenina, que se forma
sin ninguna gónada; el testículo sí hace falta para la vía masculina y
para la supresión de los conductos femeninos.

\textbf{2.} Después del día 22 los conductos ya han empezado a
diferenciarse bajo la influencia de la gónada; el experimento debe
actuar antes de la decisión.

\textbf{3.} Las hormonas del testículo, llevadas por la sangre hasta
los conductos: el injerto, lejos de los conductos, los masculiniza
igualmente.

\textbf{4.} La \omterm{prop:g11:becoming-male-female:hormones}{testosterona} construye los conductos masculinos; no
suprime los femeninos.

\textbf{5.} Una segunda sustancia testicular debe hacer regresar los
conductos de Müller; podría confirmarse extrayéndola de testículos e
implantándola sola, con la expectativa de que desaparezcan los
conductos femeninos y no se formen los masculinos --- que es lo que
hace la AMH.

\textbf{6.} En el lado operado, sin \omterm{prop:g11:becoming-male-female:hormones}{testosterona} ni AMH: conducto
femenino y ningún conducto masculino. En el lado intacto: conducto
masculino y conducto femenino regresado.

\textbf{7.} Del lado del injerto, conducto masculino y ningún conducto
femenino; del otro lado, solo conducto femenino.

\textbf{8.} Del lado del cristal, los dos conductos presentes; del otro
lado, solo el conducto femenino.

\textbf{9.} Cada feto sirve de testigo de sí mismo: la diferencia entre
los dos lados solo puede deberse a la presencia local del testículo, de
la \omterm{prop:g11:becoming-male-female:hormones}{testosterona} o a su ausencia.

\textbf{10.} Sus ovarios están extirpados o anulados, y el testículo
injertado, de otro animal, no fabrica espermatozoides en un huésped
ajeno; la vía es masculina, pero no se producen gametos.

\textbf{11.} \omterm{prop:g11:becoming-male-female:sry}{SRY} presente: testículo; \omterm{prop:g11:becoming-male-female:hormones}{testosterona} y AMH: vía masculina
y anatomía masculina; en la \omterm{prop:g11:becoming-male-female:puberty}{pubertad}, \omterm{prop:g11:becoming-male-female:hormones}{testosterona}: caracteres
masculinos. Estéril (sin los \omterm{def:g10:universal-dna:gene}{genes} del Y para los espermatozoides).

\textbf{12.} Sin un \omterm{prop:g11:becoming-male-female:sry}{SRY} funcional: no hay testículo; sin hormonas: vía
y anatomía femeninas; gónadas poco desarrolladas.

\textbf{13.} Testículos; conductos de Wolff desarrollados (la
\omterm{prop:g11:becoming-male-female:hormones}{testosterona} actúa); conductos de Müller que persisten también (no hay
respuesta a la AMH): un varón que tiene, además, útero y oviductos;
anatomía y \omterm{prop:g11:becoming-male-female:puberty}{pubertad} masculinas.

\textbf{14.} Solo el tercero: testículos funcionales y conductos
masculinos, aunque el útero que persiste puede dificultar el descenso
de los testículos. Al primero le faltan los \omterm{def:g10:universal-dna:gene}{genes} de los
espermatozoides; el segundo no tiene gónadas funcionales.

\textbf{15.} El sexo cromosómico predice el sexo anatómico solo a
través de los pasos intermedios: un \omterm{prop:g11:becoming-male-female:sry}{SRY} que funcione, hormonas que
funcionen, receptores que funcionen. Cambia cualquier paso y los dos
niveles se separan.

\textbf{16.} Chicas: el 20\% a los 10, estirón desde los 11
aproximadamente, parada hacia los 15. Chicos: el 20\% hacia los 12,5,
estirón desde los 13,5 aproximadamente, parada hacia los 17.

\textbf{17.} Las chicas, unos 4 años; los chicos, unos 3,5 --- pero los
chicos parten de una base más alta, porque han crecido al ritmo
infantil dos años más.

\textbf{18.} Los chicos, unos $3.5 \times 8 = \qty{28}{cm}$, y las
chicas $4 \times
7 = \qty{28}{cm}$: los estirones son parecidos; la
diferencia adulta viene sobre todo de los dos años más de crecimiento
infantil anteriores al estirón de los chicos.

\textbf{19.} Una subida temprana lleva los \omterm{def:g10:muscles-and-joints:joint}{cartílagos} de crecimiento al
cierre años antes: el niño crece deprisa ahora pero se detiene en una
estatura adulta baja. Retrasar las hormonas prolonga el crecimiento
infantil.

\textbf{20.} La \omterm{prop:g11:becoming-male-female:hormones}{testosterona}, que construye los conductos masculinos, y
la hormona antimülleriana, que suprime los femeninos; las dos las
fabrica el testículo, un órgano que enciende el único \omterm{prop:g11:becoming-male-female:sry}{gen SRY}.
\end{solution}
